\subsection{Interpretacja wyników}

\paragraph{Brak wariantu lepszego od wariantu odniesienia.} Rysunek~\ref{fig:sila} pokazuje kompromis między siłą a~niezależnością drzew. Przy drzewach bez ograniczenia głębokości wariant odniesienia leży na tej krzywej blisko optimum dla typowego zbioru. Zmniejszenie próbki zwiększa niezależność drzew, ale w~większości zbiorów strata siły przeważa nad tym zyskiem. Poprawa pojawia się w~zbiorach, w~których zmienna docelowa jest silnie zaszumiona (Adult, Abalone, Diabetes, Wine Quality). Pojedyncze drzewo dopasowuje się wtedy do szumu, a~mniejsza próbka działa jak regularyzacja. Jest to zgodne z~wynikiem Martíneza-Muñoza i~Suáreza~\cite{martinez2010}, że optymalny rozmiar próbki zależy od zbioru i~bywa mniejszy niż $n$. Takie zbiory stanowią mniejszość wśród naszych 24, dlatego mediana efektu mniejszych próbek jest ujemna. Z~wyników nie wynika więc, że mniejsze próbki są gorsze, tylko że rozmiar próbki należy dobierać dla zbioru, np.\ na podstawie błędu na obserwacjach spoza próbki (\emph{out-of-bag}), zamiast zmieniać go dla wszystkich zbiorów jednakowo.

\paragraph{Równoważność S1 i~S2 przy tym samym $U$.} Drzewo bez ograniczenia głębokości traktuje powtórzenia obserwacji jak wagi w~kryterium podziału. W~liściu z~jedną obserwacją waga nie zmienia predykcji, więc o~kształcie drzewa decyduje głównie zbiór różnych obserwacji. Nasz wynik potwierdza empirycznie obserwację Bühlmanna i~Yu~\cite{buhlmann2002}, że losowanie bez zwracania z~odpowiednio dobranym rozmiarem próbki działa podobnie jak losowanie bootstrapowe. Wariant S2 z~$s\approx0.63$ można więc stosować zamiennie z~wariantem odniesienia.

\paragraph{Szkodliwość pokrywania się próbek.} S3 jest jedyną interwencją w~badaniu, która zmienia niezależność próbek przy stałej sile drzew (poziomy pas na Rysunku~\ref{fig:sila}). Wynik jest jednoznaczny. Wspólny rdzeń zwiększa korelację drzew i~przy stałej sile pogarsza zespół, zgodnie z~oszacowaniem Breimana. Zwiększa też zmienność wyniku między podziałami ($\sigma/\sigma_0=1.20$ dla $\gamma=0.75$ w~regresji). Cały zespół zależy wtedy od jednego losowania rdzenia, a~uśrednianie po drzewach nie usuwa tej części losowości. W~lesie losowym efekt jest słabszy, bo losowanie cech zmniejsza korelację drzew niezależnie od próbek (Sekcja~\ref{sec:rf-bagging}).

\paragraph{Słaba zdolność predykcyjna miar próbek.} Jest to główny wynik negatywny projektu i~ma dwie przyczyny.
\begin{enumerate}
  \item \emph{W~obrębie jednego wariantu} miary są prawie stałe (Sekcja~\ref{sec:koncentracja}). Różnice między próbkami tego samego zespołu są rzędu $1/\sqrt{n}$. O~różnicach między drzewami decyduje to, \emph{które} obserwacje trafiły do próbki, a~nie \emph{ile} z~nich jest wspólnych, a~miary oparte na liczbie wspólnych obserwacji nie zawierają tej informacji. Dlatego związki w~H1 są istotne, ale bardzo słabe ($\rho\approx0.02$--$0.03$), a~w~eksperymencie dodatkowym znikają.
  \item \emph{Między wariantami} miary dobrze opisują sposób losowania, ale jakość zależy od niego w~dwóch przeciwnych kierunkach. Rozmiar próbki wpływa na siłę drzew (opisuje go $U$), a~pokrywanie się próbek na ich niezależność (opisuje je $D_J$ przy stałym $U$). Każda pojedyncza miara łączy oba kierunki, więc jej związek z~jakością jest niemonotoniczny (Rysunek~\ref{fig:miary-rangi}). Związek staje się spójny dopiero po uwzględnieniu siły drzew, co pokazują korelacje cząstkowe w~H2.
\end{enumerate}
Wynik ten jest zgodny z~obserwacją Kunchevej i~Whitaker~\cite{kuncheva2003}, że nawet miary różnorodności liczone na predykcjach nie mają prostego związku z~dokładnością zespołu. Miary próbek dobrze opisują więc \emph{sposób} losowania, czyli rozmiar próbek i~stopień ich pokrywania się, ale nie pozwalają przewidzieć, które konkretne losowanie da lepszy zespół.

\paragraph{Brak wpływu reprezentatywności.} Próbka bootstrapowa zawiera około 63\% różnych obserwacji części treningowej, więc rozkłady cech i~zmiennej docelowej w~próbce są bardzo bliskie rozkładom w~części treningowej. Losowe odchylenia $R_Y$ i~$R_X$ są niewielkie i~nie przekładają się na jakość drzewa (H1). Losowanie warstwowe, które całkowicie usuwa wahania proporcji klas, nie zmienia wyniku w~sposób praktycznie istotny w~żadnym zbiorze. Największą poprawę daje w~zbiorach z~wieloma klasami, w~tym rzadkimi (Glass $+1.3$\pp, Vehicle $+0.9$\pp), bo zapewnia obecność każdej klasy w~każdej próbce. Nawet tam poprawa nie przekracza jednak zmienności wyniku między podziałami. Wariant S4 nie pogarsza więc wyników i~może być przydatny przy wielu klasach, ale w~ogólności nie poprawia zespołu.

\paragraph{Różnice między klasyfikacją a~regresją.} Regresja jest znacznie bardziej wrażliwa na wariant losowania. S1~$r=0.2$ zwiększa w~niej RMSE w~medianie o~10\%, a~w~klasyfikacji obniża macro-F1 o~1.6\pp{}. Domyślny las losowy w~regresji nie losuje cech, więc jedynym źródłem zmniejszenia korelacji drzew jest losowanie próbek. Ponadto RMSE reaguje na każdą zmianę predykcji, a~macro-F1 tylko na zmianę przewidywanej klasy. Z~tego powodu wyniki regresji są bliższe wynikom baggingu w~klasyfikacji niż wynikom lasu losowego.

\paragraph{Ocena przyjętych rozwiązań.} Dobrze sprawdziły się trzy elementy projektu: wariant S3, który rozdziela siłę i~niezależność drzew; sparowane porównania na wspólnych podziałach i~ziarnach; analiza łączna i~test bayesowski, które pozwoliły odróżnić brak różnicy (S4, S2~$s=0.63$) od braku mocy testu (małe próbki). Nie potwierdziło się natomiast założenie, że proste miary próbek pozwolą wybrać wariant lub konkretne losowanie bez uczenia modeli. Wynika to bezpośrednio z~koncentracji miar (Sekcja~\ref{sec:koncentracja}): przy ustalonym wariancie miary prawie się nie zmieniają, więc nie mogą wyjaśniać zmienności jakości zespołu.
